\chapter[Równania półliniowe i eksplozja]{Równania półliniowe: lokalne istnienie, przedłużanie, eksplozja i energia}
\label{chap:pde25}

\section{Model, pojęcia i konieczność doboru przestrzeni}
Równanie jest \emph{półliniowe}, gdy najwyższe pochodne niewiadomej
występują liniowo, a ich współczynniki nie zależą od niewiadomej.
Przykładami są $u_t-\Delta u=f(u)$ oraz
$u_{tt}-\Delta u=f(u)$; natomiast $u_t-\diver(a(u)\nabla u)=0$
jest na ogół równaniem quasi-liniowym.
Dla zwięzłości pełny dowód lokalnego istnienia przeprowadzimy
najpierw na płaskim torusie $\T^d=\R^d/(2\pi\mathbb Z)^d$.
Przez $C(\T^d)$ rozumiemy rzeczywistą przestrzeń Banacha funkcji
ciągłych z normą supremum. W przeciwieństwie do znormalizowanej
normy $L^2(\T)$ z rozdziału~\ref{chap:pde-fourier},
tu używamy nieznormalizowanej miary objętości. Miara $\lambda_d$ na torusie jest
miarą Lebesgue'a przeniesioną z kostki $[0,2\pi)^d$;
$|\T^d|=(2\pi)^d$. Na torusie nie ma warunków brzegowych,
więc nieliniowość i mechanizmy rozwoju osobliwości nie są zasłaniane
kwestią regularności brzegu.

\begin{definicja}[Rozwiązanie łagodne]
\label{def:pde25-mild}
Niech $u_0\in C(\T^d)$, $f\in C^1(\R)$ i $T>0$.
Funkcję $u\in C([0,T];C(\T^d))$ nazywamy
\emph{rozwiązaniem łagodnym} równania $u_t-\Delta u=f(u)$
z warunkiem $u(0)=u_0$, jeżeli
\begin{equation}\label{eq:pde25-duhamel}
 u(t)=S(t)u_0+\int_0^tS(t-s)f(u(s))\,\mathrm d\lambda_1(s),
 \qquad 0\le t\le T,
\end{equation}
gdzie $S(t)=e^{t\Delta}$ jest półgrupą cieplną na $\T^d$,
a $f(u(s))$ oznacza złożenie funkcji $f$ z $u(s,\cdot)$.
Całka w~\eqref{eq:pde25-duhamel} jest całką Bochnera
w przestrzeni Banacha $C(\T^d)$.
\end{definicja}

\begin{lemat}[Półgrupa ciepła na torusie]
\label{lem:pde25-heat}
Dla każdego $d\ge1$ operatory $S(t)$, $t\ge0$, tworzą silnie
ciągłą półgrupę na $C(\T^d)$. Są dodatnie, zachowują stałe
i spełniają $\|S(t)v\|_\infty\le\|v\|_\infty$.
Ponadto $\int_{\T^d}S(t)v\,\mathrm d\lambda_d
=\int_{\T^d}v\,\mathrm d\lambda_d$.
\end{lemat}
\begin{proof}
Dla $t>0$ periodyzujemy jądro Gaussa:
$K_t(z)=\sum_{k\in\mathbb Z^d}(4\pi t)^{-d/2}
\exp(-|z+2\pi k|^2/(4t))$.
Na kostce fundamentalnej $|z+2\pi k|\ge 2\pi|k|-2\pi\sqrt d$.
Każda pochodna przestrzenna składnika szeregu jest
wielomianem w $z+2\pi k$ pomnożonym przez funkcję Gaussa.
Na zwartym przedziale dodatnich czasów szeregi tych
pochodnych mają zatem sumowalną majorantę, niezależną
od $z$ i $t$. To samo dotyczy pochodnych czasowych.
Mamy $K_t\ge0$ i
$\int_{\T^d}K_t(z)\,\mathrm d\lambda_d(z)=1$,
ponieważ kostki przesunięte o $2\pi k$ pokrywają
$\R^d$ z dokładnością do zbiorów miary zero, a całka
jądra Gaussa na $\R^d$ wynosi $1$;
stąd dodatniość, kontrakcja i zachowanie stałych dla
$S(t)v=K_t*v$. Splot jąder Gaussa na $\R^d$ oraz
periodyzacja dają $S(t)S(s)=S(t+s)$. Dla $v$ ciągłej
i jednostajnie ciągłej na torusie rozdzielamy całkę splotową
na $|z|<\delta$ i $|z|\ge\delta$ (odległość liczona na torusie).
Pierwszą część oszacowujemy przez moduł ciągłości $v$,
a druga jest ograniczona przez $2\|v\|_\infty$
pomnożone przez masę jądra poza kulą torusową
o promieniu $\delta$. Ta masa nie przekracza
$\int_{|z|\ge\delta}(4\pi t)^{-d/2}
e^{-|z|^2/(4t)}\,\mathrm d\lambda_d(z)$,
gdyż odległość torusowa od zera nie przekracza
odległości euklidesowej żadnego reprezentanta.
Po podstawieniu $z=\sqrt t\,w$ ostatnia całka
dąży do zera przez całkowalność funkcji Gaussa.
Zatem $\|S(t)v-v\|_\infty\to0$. Tożsamość całek wynika
z twierdzenia Fubiniego i jednostkowej masy jądra.
\end{proof}

\begin{uwaga}[Dlaczego nie zaczynamy od $L^2$?]
Dla wielomianowej nieliniowości $f(u)=u^3$ operator
Nemyckiego $u\mapsto f(u)$ nie jest lokalnie
lipschitzowski z $L^2(\T^d)$ do $L^2(\T^d)$.
Istotnie wybierzmy mierzalne zbiory $E_n$ o mierze
$n^{-4}$ i połóżmy $v_n=n\chi_{E_n}$.
Wówczas $\|v_n\|_2=n^{-1}\to0$, ale
$\|v_n^3\|_2=n\to\infty$.
Ponadto na małej kuli w kostce fundamentalnej
funkcja $v(x)=|x|^{-d/4}$, przedłużona przez zero
poza kulą, należy do $L^2$, lecz jej sześcian
nie należy do $L^2$ (całki radialne mają wykładniki
odpowiednio $d/2-1$ oraz $-d/2-1$). Twierdzenie Picarda w tej przestrzeni
nie byłoby więc uzasadnione. Przestrzeń $C(\T^d)$ pozwala
przeprowadzić argument bez założeń o wykładnikach Sobolewa.
\end{uwaga}

\section{Twierdzenie Picarda i kryterium przedłużania}
\begin{twierdzenie}[Lokalne istnienie i jednoznaczność]
\label{thm:pde25-local}
Niech $d\ge1$, $u_0\in C(\T^d)$ i $f\in C^1(\R)$.
Istnieje $T_0>0$ i dokładnie jedno rozwiązanie łagodne
$u\in C([0,T_0];C(\T^d))$ równania~\eqref{eq:pde25-duhamel}.
Rozwiązanie posiada jednoznaczne maksymalne przedłużenie
na przedział $[0,T_{\max})$, gdzie $T_{\max}\in(0,\infty]$.
Jeżeli $T_{\max}<\infty$, to
\begin{equation}\label{eq:pde25-blowalt}
 \limsup_{t\uparrow T_{\max}}\|u(t)\|_\infty=\infty.
\end{equation}
W rzeczywistości w takim przypadku norma dąży do
nieskończoności w sensie granicy.
\end{twierdzenie}
\begin{proof}
\textbf{Krok 1: wybór kuli.} Ustalmy $R>\|u_0\|_\infty+1$.
Na $[-R,R]$ przyjmijmy
$L_R=\max_{|z|\le R}|f'(z)|$ i
$M_R=\max_{|z|\le R}|f(z)|$.
W zupełnej przestrzeni $Z_T=C([0,T];C(\T^d))$
rozważmy domkniętą kulę
$B_R=\{v:\|v\|_{Z_T}\le R\}$.
Definiujemy $\Phi(v)$ jako prawą stronę~\eqref{eq:pde25-duhamel}.
Dla $v\in Z_T$ superpozycja $s\mapsto f(v(s))$ jest
ciągła w normie supremum, gdyż $f$ jest lipschitzowska
na ograniczonym zakresie wartości $v$.
Silna ciągłość i kontraktywność półgrupy gwarantują
ciągłość całki Bochnera względem górnej granicy:
dla $t_n\to t$ rozpatrujemy całki na wspólnym
przedziale i stosujemy zbieżność zmajoryzowaną w normie,
a całkę po różnicy przedziałów ograniczamy przez
$|t_n-t|\sup_{s\in[0,T]}\|f(v(s))\|_\infty$
(dla $v\in B_R$ jest to co najwyżej $|t_n-t|M_R$). Zatem $\Phi$ rzeczywiście przyjmuje
wartości w $Z_T$.
Lemat~\ref{lem:pde25-heat} daje
$\|\Phi(v)(t)\|_\infty\le\|u_0\|_\infty+TM_R$.
Wybierając $T$ tak, by $TM_R\le R-\|u_0\|_\infty$,
otrzymujemy $\Phi(B_R)\subseteq B_R$.

\textbf{Krok 2: kontrakcja.} Twierdzenie o wartości średniej
i kontrakcja półgrupy implikują
$\|\Phi(v)-\Phi(w)\|_{Z_T}\le TL_R\|v-w\|_{Z_T}$.
Jeśli $TL_R<1$, zasada Banacha daje jedyny punkt stały
w $B_R$. Jednoznaczność nie zależy od wybranej kuli:
dla dowolnych dwóch rozwiązań na wspólnym przedziale
ich zakresy są ograniczone przez pewne $R$, a nierówność
$\|u(t)-v(t)\|_\infty\le
L_R\int_0^t\|u(s)-v(s)\|_\infty\,\mathrm d\lambda_1(s)$
i lemat Gronwalla wymuszają równość.

\textbf{Krok 3: sklejanie.} Półgrupowość pozwala przepisać
wzór Duhamela od dowolnego czasu $t_0$:
$u(t_0+t)=S(t)u(t_0)+\int_0^tS(t-s)f(u(t_0+s))\,\mathrm d\lambda_1(s)$.
Tożsamość otrzymujemy, rozdzielając całkę w chwili
$t_0$ i stosując własność półgrupy do pierwszej części.
Odwrotnie, rozwiązanie rozpoczynające się od $u(t_0)$
skleja się z poprzednim, gdyż po podstawieniu wzoru
na $u(t_0)$ spełnia pierwotny wzór Duhamela.
Niech $T_{\max}$ będzie supremum czasów istnienia.
Jednoznaczność na wspólnych przedziałach pozwala
zdefiniować rozwiązanie na $[0,T_{\max})$ przez
zgodne ograniczenia; każde jego przedłużenie poza
ten przedział przeczyłoby definicji supremum.

\textbf{Krok 4: alternatywa eksplozji.} Jeśli $T_{\max}<\infty$
i $\|u(t)\|_\infty\le R$ dla wszystkich $t$ bliskich
$T_{\max}$, to w konstrukcji z kroku pierwszego można
zastosować stale tę samą większą kulę oraz ten sam czas
lokalnego istnienia $\tau_R>0$, niezależny od chwili restartu.
Wybierając $t_0\in(T_{\max}-\tau_R/2,T_{\max})$
przedłużamy rozwiązanie poza $T_{\max}$, sprzeczność.
Ten sam argument pokazuje, że nie może istnieć ciąg
$t_n\uparrow T_{\max}$ z $\|u(t_n)\|_\infty\le R$,
ponieważ restart w $t_n$ daje przedłużenie o stały czas.
Zatem norma rzeczywiście dąży do nieskończoności.
\end{proof}

\begin{twierdzenie}[Stabilność i ciągła zależność od danych]
\label{thm:pde25-stability}
Niech $u,v$ będą rozwiązaniami łagodnymi na $[0,T]$,
z początkami $u_0,v_0\in C(\T^d)$, dla tej samej
$f\in C^1(\R)$. Jeśli
$\max\{\|u\|_{Z_T},\|v\|_{Z_T}\}\le R$, to
$\|u(t)-v(t)\|_\infty\le e^{L_Rt}\|u_0-v_0\|_\infty$
dla każdego $0\le t\le T$.
\end{twierdzenie}
\begin{proof}
Odejmujemy wzory Duhamela, używamy oszacowania
$|f(a)-f(b)|\le L_R|a-b|$ dla $|a|,|b|\le R$,
a następnie lematu Gronwalla w postaci całkowej.
\end{proof}


\begin{uwaga}[Przenoszenie konstrukcji między rozdziałami]
\label{rem:pde-v44-picard-scaling}
Dowód kontrakcji z twierdzenia~\ref{thm:pde25-local} opiera się na
silnej ciągłości i kontraktywności półgrupy w $C(\T^d)$ oraz na
lokalnej lipschitzowskości operatora Nemyckiego w tej normie.
W rozdziale~\ref{chap:pde27} użyjemy dokładnie tego samego
mechanizmu na torusie o okresie $L$ i z dyfuzją $D>0$;
zmiana okresu i czasu zostanie tam uzasadniona wzorem
\eqref{eq:pde-v44-torus-scale}. Nie jest natomiast poprawne
przeniesienie dowodu bezpośrednio na $L^2$ dla dowolnej
nieliniowości wielomianowej: odwzorowanie $v\mapsto v^3$
nie musi przyjmować wartości w $L^2$ ani być lokalnie
lipschitzowskie w tej przestrzeni.
\end{uwaga}


\section{Zasada porównania, dodatniość i istnienie globalne}
\begin{twierdzenie}[Porównanie rozwiązań półliniowych]
\label{thm:pde25-comparison}
Niech $f\in C^1(\R)$, $T>0$, a $u,v\in Z_T$
będą rozwiązaniami łagodnymi z danymi $u_0,v_0$.
Jeżeli $u_0(x)\le v_0(x)$ dla każdego $x\in\T^d$,
to $u(t,x)\le v(t,x)$ dla wszystkich
$(t,x)\in[0,T]\times\T^d$.
Nie zakładamy monotoniczności $f$.
\end{twierdzenie}
\begin{proof}
Oznaczmy $w=v-u$ i
$a(t,x)=\int_0^1f'(u+\theta(v-u))\,\mathrm d\lambda_1(\theta)$.
Wtedy $f(v)-f(u)=aw$ i $|a|\le L$ dla pewnej
stałej $L$ zależnej od zakresów rozwiązań.
Dla $z(t,x)=e^{Lt}w(t,x)$ otrzymujemy równanie
\[
 z(t)=e^{Lt}S(t)w(0)+
 \int_0^t e^{L(t-s)}S(t-s)a(s)z(s)\,\mathrm d\lambda_1(s).
\]
Aby zmienić je na wzór z półgrupą $S$, używamy
tożsamości
\[
 e^{Lt}S(t)=S(t)+L\int_0^t S(t-r)e^{Lr}S(r)\,
 \mathrm d\lambda_1(r),
\]
wynikającej z własności półgrupy i całki z $Le^{Lr}$.
Podstawiając ją do obu składników wzoru dla $z$
oraz zamieniając kolejność całkowania w trójkącie
$0\le s\le r\le t$, otrzymujemy
\[
 z(t)=S(t)w(0)+
 \int_0^t S(t-s)((a(s)+L)z(s))\,\mathrm d\lambda_1(s).
\]
Zamiana kolejności jest dopuszczalna dla całek
Bochnera dzięki ciągłości i ograniczoności
funkcji podcałkowych na zwartym trójkącie.
Ten sam rachunek działa dla dodatkowego źródła
$b(t)$ w równaniu dla $w$: w ostatnim wzorze
pojawia się dodatkowo $e^{Ls}b(s)$ pod całką.
Tu $a+L\ge0$. Kolejne iteracje tego liniowego równania
Volterry, zaczynając od $z_0(t)=S(t)w(0)\ge0$, pozostają
nieujemne, gdyż półgrupa jest dodatnia. Iteracje zbiegają
jednostajnie na dostatecznie krótkim przedziale przez
kontrakcję ($T\|a+L\|_\infty<1$). Na pozostałych
częściach $[0,T]$ powtarzamy ten sam argument.
Otrzymujemy $w\ge0$.
\end{proof}

\begin{wniosek}[Globalność przy wzroście liniowym]
\label{cor:pde25-global}
Jeśli $f\in C^1(\R)$ i dla pewnych $A,B\ge0$ zachodzi
$|f(s)|\le A+B|s|$ dla każdego $s\in\R$, to każde
rozwiązanie łagodne na $\T^d$ istnieje dla wszystkich
$t\ge0$. Ponadto
\[
 \|u(t)\|_\infty\le
 e^{Bt}\|u_0\|_\infty+\frac AB(e^{Bt}-1)
\]
dla $B>0$; przy $B=0$ ostatni składnik należy zastąpić
przez $At$.
\end{wniosek}
\begin{proof}
Z kontrakcji półgrupy i wzoru Duhamela wynika
$\|u(t)\|_\infty\le\|u_0\|_\infty+
\int_0^t(A+B\|u(s)\|_\infty)\,\mathrm d\lambda_1(s)$.
Gronwall daje żądane ograniczenie na każdym skończonym
przedziale. Alternatywa~\eqref{eq:pde25-blowalt}
wyklucza skończony czas maksymalnego istnienia.
\end{proof}

\begin{przyklad}[Równanie logistyczne]
\label{ex:pde25-logistic}
Dla $u_t=\Delta u+u(1-u)$ i
$0\le u_0\in C(\T^d)$ otrzymujemy
$0\le u(t,x)\le M=\max\{1,\|u_0\|_\infty\}$.
Funkcja stała $0$ jest rozwiązaniem, a $M$ jest
supersolucją, bo $M(1-M)\le0$. Porównanie z funkcją
stałą $M$ można uzasadnić bez dodatkowej teorii
supersolucji: niech $m(t)$ rozwiązuje
$m'=m(1-m)$, $m(0)=M$; wtedy $m(t)\le M$.
Twierdzenie~\ref{thm:pde25-comparison} zastosowane
do rozwiązań $u$ i $m$ daje ograniczenie i globalność.
Choć $f(s)=s(1-s)$ rośnie kwadratowo w module,
nie występuje eksplozja dla takich danych nieujemnych.
\end{przyklad}


\begin{twierdzenie}[Oszacowanie przez równanie skalarne; warunek Osgooda]
\label{thm:pde25-osgood}
Niech $d\ge1$, $u_0\in C(\T^d)$ i $f\in C^1(\R)$.
Załóżmy, że funkcja $h:[0,\infty)\to(0,\infty)$ jest
niemalejąca, lokalnie lipschitzowska i spełnia
$|f(s)|\le h(|s|)$ dla wszystkich $s\in\R$.
Niech $M=\|u_0\|_\infty$, a $y:[0,\tau_M)\to[0,\infty)$
będzie maksymalnym rozwiązaniem $y'=h(y)$, $y(0)=M$.
Wówczas maksymalne rozwiązanie łagodne równania
$u_t=\Delta u+f(u)$ spełnia
$T_{\max}\ge\tau_M$ oraz
$\|u(t)\|_\infty\le y(t)$ dla $0\le t<\tau_M$.
Ponadto
\begin{equation}\label{eq:pde25-osgood}
 \tau_M=\int_M^\infty\frac{\mathrm d\lambda_1(s)}{h(s)},
\end{equation}
gdzie dopuszczamy wartość $+\infty$.
W szczególności rozbieżność całki w~\eqref{eq:pde25-osgood}
gwarantuje globalne istnienie rozwiązania $u$.
\end{twierdzenie}
\begin{proof}
Funkcja $H(r)=\int_M^r h(s)^{-1}\,\mathrm d\lambda_1(s)$,
$r\ge M$, jest klasy $C^1$ i ma ściśle dodatnią
pochodną. Jej obrazem jest przedział $[0,\tau_M)$,
gdzie $\tau_M$ jest prawą stroną~\eqref{eq:pde25-osgood}.
Odwrotność $y=H^{-1}$ spełnia $y(0)=M$ oraz
$y'=1/H'(y)=h(y)$. Każde rozwiązanie ma tę postać
przez regułę łańcuchową stosowaną do $H(y(t))$.
Jeżeli $\tau_M<\infty$, to $y(t)\to\infty$ przy
$t\uparrow\tau_M$, więc przedłużenie nie istnieje.
Jeżeli $\tau_M=\infty$, odwrotność jest określona
dla wszystkich nieujemnych czasów.

Ustalmy $T<\min\{T_{\max},\tau_M\}$. Funkcje stałe
w przestrzeni $\overline y(t,x)=y(t)$ oraz
$\underline y(t,x)=-y(t)$ spełniają odpowiednio
$\overline y_t-\Delta\overline y-f(\overline y)\ge0$
i $\underline y_t-\Delta\underline y-f(\underline y)\le0$.
Dla uzasadnienia porównania w klasie łagodnej zauważmy,
że $w=\overline y-u$ spełnia wzór Duhamela z nieujemnym
danym początkowym $M-u_0$ oraz dodatkowym nieujemnym
źródłem $h(y)-f(y)$. Pozostała część źródła jest równa
$a(t,x)w$, gdzie
$a=\int_0^1f'(u+\theta(\overline y-u))\,\mathrm d\lambda_1(\theta)$
jest ograniczona. Mnożnik $e^{Lt}$, $L\ge\|a\|_\infty$,
przekształca to równanie w równanie Volterry z dodatnim
jądrem i nieujemnym współczynnikiem. Iteracja jak w dowodzie
twierdzenia~\ref{thm:pde25-comparison} daje $w\ge0$.
Analogicznie $u-\underline y\ge0$ (dodatkowe źródło wynosi
$f(-y)+h(y)\ge0$). Zatem $|u(t,x)|\le y(t)$.

Gdyby $T_{\max}<\tau_M$, funkcja $y$ byłaby ograniczona
na $[0,T_{\max}]$. Uzyskane oszacowanie przeczyłoby
alternatywie eksplozji z twierdzenia~\ref{thm:pde25-local}.
Wobec tego $T_{\max}\ge\tau_M$.
\end{proof}

\begin{przyklad}[Wzrost szybszy niż liniowy bez eksplozji z tego kryterium]
Dla $h(s)=(1+s)\log(e+s)$ całka
$\int_M^\infty h(s)^{-1}\,\mathrm d\lambda_1(s)$ jest rozbieżna,
gdyż dla $s\ge e$ zachodzi $1+s\le2s$ i
$\log(e+s)\le\log(2s)\le2\log s$, a całka
$\int_e^\infty (s\log s)^{-1}\,\mathrm d\lambda_1(s)$
jest rozbieżna po podstawieniu $r=\log s$.
Każda nieliniowość $f\in C^1(\R)$ spełniająca
$|f(s)|\le(1+|s|)\log(e+|s|)$ daje więc globalne
rozwiązania na torusie dla dowolnych danych ciągłych.
Pokazuje to, że warunek wzrostu liniowego z
wniosku~\ref{cor:pde25-global} nie jest konieczny.
\end{przyklad}


\section{Eksplozja w skończonym czasie: zasada średniej}
\begin{twierdzenie}[Wybuch dla dodatniego źródła na torusie]
\label{thm:pde25-blowup}
Niech $d\ge1$, $p>1$ oraz
$u_0\in C(\T^d)$, $u_0\ge0$, $u_0\not\equiv0$.
Niech $u$ będzie maksymalnym rozwiązaniem łagodnym
równania $u_t=\Delta u+(u_+)^p$, gdzie $u_+=\max\{u,0\}$.
Oznaczmy $m_0=|\T^d|^{-1}\int_{\T^d}u_0(x)
\,\mathrm d\lambda_d(x)>0$. Wówczas $u\ge0$ i
\begin{equation}\label{eq:pde25-mass-blowup}
 T_{\max}\le\frac{m_0^{1-p}}{p-1}<\infty.
\end{equation}
\end{twierdzenie}
\begin{proof}
Funkcja $s\mapsto(s_+)^p$ należy do $C^1(\R)$,
więc twierdzenie lokalne obowiązuje. Porównanie z
rozwiązaniem zerowym daje $u\ge0$. Z zachowania masy
przez półgrupę i~\eqref{eq:pde25-duhamel} wynika
\[
 m(t)=m_0+\int_0^t
 \frac1{|\T^d|}\int_{\T^d}u(s,x)^p\,
 \mathrm d\lambda_d(x)\,\mathrm d\lambda_1(s).
\]
Funkcja $m$ jest klasy $C^1$ i, na mocy nierówności
Jensena dla wypukłej funkcji $s\mapsto s^p$ na
$[0,\infty)$, spełnia $m'(t)\ge m(t)^p$.
Ponieważ $m(t)>0$, możemy podzielić przez $m(t)^p$:
$\frac{\mathrm d}{\mathrm dt}m(t)^{1-p}
=(1-p)m(t)^{-p}m'(t)\le-(p-1)$.
Stąd $m(t)^{1-p}\le m_0^{1-p}-(p-1)t$.
Lewa strona jest dodatnia na całym przedziale istnienia,
co wymusza~\eqref{eq:pde25-mass-blowup}.
\end{proof}

\begin{uwaga}[Źródło eksplozji]
Wzrost normy wynika tu z nierówności różniczkowej dla
średniej, a nie z nieciągłości współczynników lub
braku regularności danych. Co więcej, nawet niewielkie,
niezerowe dane nieujemne prowadzą do eksplozji na torusie.
To twierdzenie nie przenosi się bez zmian na $\R^d$,
gdzie rozpraszanie masy może zapobiegać wybuchowi.
Wymagane są tam inne argumenty i krytyczne wykładniki.
\end{uwaga}

\begin{twierdzenie}[Wybuch dla dużych danych Dirichleta]
\label{thm:pde25-kaplan}
Niech $\Omega\subset\R^d$ będzie ograniczonym obszarem
klasy $C^2$, $p>1$ i $u$ nieujemnym rozwiązaniem klasycznym
na maksymalnym przedziale $[0,T_{\max})$
z regularnością klasyczną aż do brzegu zagadnienia
$u_t-\Delta u=u^p$, $u|_{\partial\Omega}=0$.
Niech $\lambda_1>0$ będzie pierwszą wartością własną
$-\Delta$ z warunkiem Dirichleta, a $\phi_1>0$
odpowiadającą jej funkcją własną
(Twierdzenie~\ref{thm:pde18-simple}),
$\int_\Omega\phi_1\,\mathrm d\lambda_d=1$.
Jeśli
$y_0=\int_\Omega u_0(x)\phi_1(x)\,\mathrm d\lambda_d(x)
>\lambda_1^{1/(p-1)}$,
to $T_{\max}<\infty$ i
$\|u(t)\|_\infty\to\infty$ dla $t\uparrow T_{\max}$,
przy czym wniosek o granicy normy wymaga założenia,
że w rozważanej klasie istnieje jednostajny dodatni
czas lokalnego przedłużenia dla danych o ograniczonej
normie supremum. Sam wniosek $T_{\max}<\infty$
nie wymaga tego dodatkowego założenia.
\end{twierdzenie}
\begin{proof}
Dla $y(t)=\int_\Omega u(t,x)\phi_1(x)
\,\mathrm d\lambda_d(x)$ korzystamy ze słabego równania
własnego dla $\phi_1\in H_0^1(\Omega)$ i z warunku
$u(t)\in H_0^1(\Omega)$:
\[
 \int_\Omega\Delta u\,\phi_1\,\mathrm d\lambda_d
 =-\int_\Omega\nabla u\cdot\nabla\phi_1\,\mathrm d\lambda_d
 =-\lambda_1\int_\Omega u\phi_1\,\mathrm d\lambda_d .
\]
Różniczkowanie pod całką i równanie paraboliczne dają
$y'=-\lambda_1y+\int_\Omega u^p\phi_1\,
\mathrm d\lambda_d$.
Miara $\phi_1\,\mathrm d\lambda_d$ ma masę $1$,
więc Jensen zapewnia $y'\ge-\lambda_1y+y^p$.
Ponieważ $y_0>\lambda_1^{1/(p-1)}$, prawa strona
jest dodatnia w otoczeniu $y_0$. Funkcja $y$ nie może
spaść poniżej $y_0$: przy pierwszym powrocie do
tej wartości z góry jej pochodna byłaby niedodatnia,
podczas gdy nierówność różniczkowa daje pochodną
ściśle dodatnią. Stąd $y$ jest ściśle rosnąca.
Ustalmy $K>\max\{y_0,(2\lambda_1)^{1/(p-1)}\}$.
Jeżeli rozwiązanie przestaje istnieć przed osiągnięciem
$K$, jego maksymalny czas już jest skończony.
W przeciwnym razie, dopóki $y\in[y_0,K]$, pochodna
jest ograniczona od dołu przez dodatnią stałą
$c=\min_{s\in[y_0,K]}(s^p-\lambda_1s)$.
Zatem przy założeniu globalnego istnienia $y$ musi
osiągnąć $K$ najpóźniej w chwili $(K-y_0)/c$.
Od tej chwili $y'\ge\frac12 y^p$, co po całkowaniu
nierówności dla $y^{1-p}$ prowadzi do sprzeczności
najpóźniej po kolejnych $2K^{1-p}/(p-1)$ jednostkach
czasu. Nie istnieje więc globalne rozwiązanie,
nawet nieograniczone w czasie.
Dokładniej, ponieważ $y'\ge y^p-\lambda_1y>0$,
całkowanie daje
\[
 t\le\int_{y_0}^{y(t)}\frac{\mathrm d\lambda_1(s)}
 {s^p-\lambda_1s}
 <\int_{y_0}^{\infty}\frac{\mathrm d\lambda_1(s)}
 {s^p-\lambda_1s}<\infty.
\]
Ostatnia całka jest skończona, gdyż mianownik
jest dodatni na $[y_0,\infty)$ i dla dużych $s$
jest co najmniej równy $s^p/2$.
Jeżeli obowiązuje założone kryterium przedłużania
z jednostajnym czasem lokalnego istnienia dla
ograniczonych danych, to argument restartu
z Twierdzenia~\ref{thm:pde25-local} daje
$\|u(t)\|_\infty\to\infty$ przy $t\uparrow T_{\max}$.
\end{proof}

\begin{twierdzenie}[Regularność rozwiązania łagodnego dla dodatnich czasów]
\label{thm:pde25-parabolic-regularity}
Niech $f\in C^1(\R)$ i niech $u$ będzie rozwiązaniem
łagodnym z Twierdzenia~\ref{thm:pde25-local}.
Dla $0<a<T<T_{\max}$ zachodzi
\[
 u\in C([a,T];C^2(\T^d))\cap C^1([a,T];C(\T^d)),
 \qquad u_t-\Delta u=f(u).
\]
Jeżeli ponadto $u_0\in H^1(\T^d)$, to
$u\in C([0,T];H^1(\T^d))$.
\end{twierdzenie}
\begin{proof}
Dla periodyzowanego jądra cieplnego z Lematu
~\ref{lem:pde25-heat} zachodzi dla $t>0$
\[
 \|\partial_j K_t\|_{L^1(\T^d)}
 \le\|\partial_j G_t\|_{L^1(\R^d)}
 =c_d t^{-1/2},
\]
gdzie $G_t$ jest jądrem Gaussa na $\R^d$.
Nierówność wynika z nierówności trójkąta i sumowania
po przesuniętych kostkach; równość z podstawienia
$z=\sqrt t\,w$ w całce z modułu pochodnej Gaussa.
Stąd $\|\partial_j S(t)v\|_\infty
\le c_dt^{-1/2}\|v\|_\infty$.

Ustalmy $0<r<a<T<T_{\max}$. Ze wzoru Duhamela
i ograniczoności $f(u)$ otrzymujemy na $[r,T]$
\[
 \partial_j u(t)=\partial_j S(t)u_0+
 \int_0^t\partial_j S(t-s)f(u(s))\,
 \mathrm d\lambda_1(s).
\]
Aby uzasadnić różniczkowanie, najpierw obcinamy
całkę do $[0,t-\delta]$. Zarówno całki, jak i ich
pochodne przestrzenne zbiegają jednostajnie na
$[r,T]\times\T^d$, gdyż ogon całki z pochodną ma
normę co najwyżej $2c_d\sqrt\delta
\sup_s\|f(u(s))\|_\infty$. Całki obcięte są ciągłe
w czasie i przestrzeni. Twierdzenie o jednostajnej
zbieżności funkcji i pochodnych daje zatem
$u\in C([r,T];C^1(\T^d))$.

Teraz rozpoczynamy wzór Duhamela w chwili $r$.
Ponieważ $\partial_j f(u)=f'(u)\partial_j u$
jest ciągła w czasie i przestrzeni na $[r,T]$,
możemy przenieść jedną pochodną na funkcję
pod splotem (przez periodyczne całkowanie przez części).
Dostajemy dla $t\ge a$
\[
 \partial_i\partial_j u(t)=\partial_i\partial_j S(t-r)u(r)
 +\int_r^t\partial_i S(t-s)
       \bigl(f'(u(s))\partial_j u(s)\bigr)\,
       \mathrm d\lambda_1(s).
\]
To samo oszacowanie ogona i jednostajna zbieżność
pochodnych uzasadniają wzór oraz ciągłość drugich
pochodnych. W szczególności $\Delta u+f(u)$ jest
ciągła na $[a,T]\times\T^d$.

Wzór Duhamela spełnia równanie dystrybucyjne:
w każdym modzie Fouriera $k\in\mathbb Z^d$ otrzymujemy
\[
 \widehat u_k'(t)=-|k|^2\widehat u_k(t)
                  +\widehat{f(u)}_k(t).
\]
Wynika to z mnożnika $e^{-t|k|^2}$ półgrupy
(otrzymanego przez całkowanie periodyzowanego jądra)
i twierdzenia Newtona--Leibniza. Całkując ten wzór
od $a$ do $t$, otrzymujemy równość wszystkich
współczynników Fouriera obu stron tożsamości
\[
 u(t)=u(a)+\int_a^t(\Delta u(s)+f(u(s)))\,
                      \mathrm d\lambda_1(s).
\]
Jednoznaczność współczynników Fouriera dla funkcji
ciągłych wynika z gęstości wielomianów trygonometrycznych.
Tożsamość w $C(\T^d)$ daje żądaną ciągłą pochodną
czasową.

Na koniec, dla $u_0\in H^1$, półgrupa jest silnie
ciągła w $H^1$ (komutuje ze słabymi pochodnymi
i jest silnie ciągła w $L^2$, co wynika z ciągłości
translacji i aproksymacji Gaussa).
Nierówność Younga i oszacowanie pochodnych jądra
dają
\[
 \left\|\int_0^t S(t-s)f(u(s))\,
              \mathrm d\lambda_1(s)\right\|_{H^1}
 \le C_d(t+\sqrt t)\sup_{0\le s\le t}\|f(u(s))\|_2.
\]
Całka Bochnera w $H^1$ istnieje dzięki tej całkowalnej
majorancie. Prawa strona dąży do zera przy $t\downarrow0$,
co wraz z ciągłością dla dodatnich czasów dowodzi
ciągłości w $H^1$ na $[0,T]$.
\end{proof}

\section{Funkcjonał energii i równanie Allena--Cahna}
\begin{twierdzenie}[Dyssypacja energii]
\label{thm:pde25-energy}
Niech $V\in C^2(\R)$ i niech $u$ będzie rozwiązaniem
łagodnym $u_t=\Delta u-V'(u)$ na $[0,T]$, z danymi
$u_0\in C(\T^d)$. Dla dodatnich czasów, zgodnie
z Twierdzeniem~\ref{thm:pde25-parabolic-regularity},
rozwiązanie jest klasyczne.
Definiujemy
\[
 E(u)=\int_{\T^d}\left(
 \frac12|\nabla u(x)|^2+V(u(x))\right)
 \,\mathrm d\lambda_d(x).
\]
Wtedy $t\mapsto E(u(t))$ jest różniczkowalna dla
$t>0$ i
\begin{equation}\label{eq:pde25-energy}
 \frac{\mathrm d}{\mathrm dt}E(u(t))
 =-\int_{\T^d}|u_t(t,x)|^2\,\mathrm d\lambda_d(x)\le0.
\end{equation}
Jeżeli dodatkowo $u_0\in H^1(\T^d)$, to dla
$0\le t\le T$ zachodzi
\[
 E(u(t))+\int_0^t\|u_t(s)\|_2^2\,
                  \mathrm d\lambda_1(s)=E(u_0).
\]
\end{twierdzenie}
\begin{proof}
Nie zakładamy istnienia mieszanej pochodnej
$\partial_t\nabla u$. Dla $t,t+h>0$ z tożsamości
różnicy kwadratów i periodycznego całkowania przez
części otrzymujemy
\[
 \begin{aligned}
 &\frac{1}{2h}\bigl(\|\nabla u(t+h)\|_2^2-
                       \|\nabla u(t)\|_2^2\bigr)\\
 &\quad=-\frac12\int_{\T^d}
    (\Delta u(t+h)+\Delta u(t))
    \frac{u(t+h)-u(t)}h\,\mathrm d\lambda_d .
 \end{aligned}
\]
Ciągłość $\Delta u$ i różniczkowalność $u$ w normie
supremum dają granicę $-\int_{\T^d}\Delta u\,u_t\,
\mathrm d\lambda_d$. Dla części potencjalnej
twierdzenie o wartości średniej w postaci całkowej
i jednostajna ciągłość $V'$ na wspólnym ograniczonym
zakresie wartości dają pochodną
$\int_{\T^d}V'(u)u_t\,\mathrm d\lambda_d$.
Z równania wynika $-\Delta u+V'(u)=-u_t$,
co dowodzi~\eqref{eq:pde25-energy}.
Pochodna jest ciągła dla dodatnich czasów.
Całkujemy od $a>0$ do $t$. Jeśli $u_0\in H^1$,
to $u(a)\to u_0$ w $H^1$ oraz w normie supremum
przy $a\downarrow0$, zatem $E(u(a))\to E(u_0)$.
Twierdzenie o zbieżności monotonicznej dla całki
z nieujemnej funkcji $\|u_t\|_2^2$ daje tożsamość
od chwili $0$, a także skończoność tej całki.
\end{proof}

\begin{przyklad}[Równanie Allena--Cahna]
Dla $V(s)=\frac14(s^2-1)^2$ równanie
$u_t=\Delta u+u-u^3$ ma stabilne, stałe stany $u\equiv\pm1$
i niestabilny stan $u\equiv0$ w sensie liniaryzacji
jednorodnych perturbacji. Istotnie dla perturbacji jednorodnej liniaryzacja
wokół stanu $a\in\{-1,0,1\}$ daje równanie
$h'=(1-3a^2)h$, a więc zanik dla $a=\pm1$
i wzrost dla $a=0$. Niech
$M\ge\max\{1,\|u_0\|_\infty\}$.
Rozwiązanie skalarne $m'=m-m^3$, $m(0)=M$,
pozostaje w $[1,M]$: pochodna jest niedodatnia
dla $m\ge1$, a jednoznaczność wyklucza przecięcie
rozwiązania stałego $1$. Jest zatem globalne.
Funkcje $m(t)$ i $-m(t)$, stałe względem $x$,
są rzeczywistymi rozwiązaniami równania półliniowego
(nie utożsamiamy ich ze stałymi $\pm M$).
Twierdzenie~\ref{thm:pde25-comparison} daje
$-M\le-m(t)\le u(t,x)\le m(t)\le M$,
a alternatywa eksplozji zapewnia globalność.
Dla dodatnich czasów energia maleje zgodnie
z~\eqref{eq:pde25-energy}; dla danych należących
również do $H^1$ zachodzi tożsamość energetyczna
od chwili początkowej.
Sama monotoniczność energii nie oznacza jednak,
że rozwiązanie musi zbiegać do konkretnego stanu
stałego bez dalszych argumentów zwartości.
\end{przyklad}

\begin{center}
\begin{tikzpicture}[x=1.65cm,y=3.0cm,>=Latex]
 \draw[->] (-1.65,0)--(1.7,0) node[right] {$s$};
 \draw[->] (0,-.045)--(0,.66) node[above] {$V(s)$};
 \draw[blue!75!black,very thick,domain=-1.52:1.52,samples=121]
  plot (\x,{.25*(\x*\x-1)^2});
 \foreach \a in {-1,1}{
  \fill[blue!75!black] (\a,0) circle (1.7pt);
  \node[below] at (\a,0) {$\a$};
 }
 \fill[red!75!black] (0,.25) circle (1.7pt);
 \draw[densely dashed,gray] (-.95,.25)--(0,.25);
 \node[left,font=\small] at (-.95,.25) {$1/4$};
 \node[below left] at (0,0) {$0$};
 \begin{scope}[xshift=6.9cm,y=1cm]
  \draw[->] (-1.65,0)--(1.7,0) node[right] {$s$};
  \foreach \a in {-1,0,1}{
   \fill (\a,0) circle (1.7pt);
   \node[below] at (\a,0) {$\a$};
  }
  \draw[->,thick,blue!75!black] (-1.55,.18)--(-1.12,.18);
  \draw[->,thick,blue!75!black] (-.15,.18)--(-.85,.18);
  \draw[->,thick,blue!75!black] (.15,.18)--(.85,.18);
  \draw[->,thick,blue!75!black] (1.55,.18)--(1.12,.18);
  \node[font=\small] at (0,.65) {$s'=s-s^3$};
 \end{scope}
\end{tikzpicture}

\small Rysunek: potencjał Allena--Cahna
$V(s)=(s^2-1)^2/4$ i kierunki ruchu dla jednorodnego
przestrzennie równania $s'=-V'(s)$. Minima potencjału
$s=\pm1$ są przyciągające, a maksimum lokalne $s=0$
jest odpychające dla tego równania skalarnego.
\end{center}

\section{Półliniowe równanie falowe w jednym wymiarze}
Dotychczas wykorzystywaliśmy wygładzającą półgrupę ciepła.
Dla równania falowego nieliniowość nie jest wygładzana:
wymagana jest przestrzeń fazowa, w której składanie
z $g$ jest lokalnie lipschitzowskim odwzorowaniem
z przestrzeni pierwszej składowej do przestrzeni
drugiej składowej. Pracujemy na $\T^1$.
Niech $A=1-\partial_x^2$ z dziedziną $H^2(\T^1)$
w przestrzeni Hilberta $L^2(\T^1)$, a
$Y=H^1(\T^1)\times L^2(\T^1)$ z normą
$\|(u,v)\|_Y^2=\|A^{1/2}u\|_2^2+\|v\|_2^2$.

\begin{twierdzenie}[Lokalne istnienie dla równania falowego]
\label{thm:pde25-wave}
Niech $g\in C^1(\R)$, $(u_0,u_1)\in Y$
i $T>0$ będzie wystarczająco mały, zależnie od normy
danych i $g$. Wtedy zagadnienie
\[
 u_{tt}-u_{xx}+u=g(u),\quad
 u(0)=u_0,\quad u_t(0)=u_1
\]
ma dokładnie jedno rozwiązanie łagodne
$(u,u_t)\in C([0,T];Y)$. Rozwiązanie zależy
lokalnie lipschitzowsko od danych w normie $Y$.
Ma jednoznaczne maksymalne przedłużenie na
$[0,T_{\max})$. Jeżeli $T_{\max}<\infty$, to
\[
 \lim_{t\uparrow T_{\max}}\|(u(t),u_t(t))\|_Y=\infty.
\]
\end{twierdzenie}
\begin{proof}
\textbf{Krok 1: grupa liniowa.} W bazie Fouriera operator
$A$ ma wartości własne $1+k^2$, $k\in\mathbb Z$.
Definiujemy $C(t)=\cos(tA^{1/2})$ i
$D(t)=A^{-1/2}\sin(tA^{1/2})$ przez mnożniki Fouriera.
Równanie liniowe ma postać
$u(t)=C(t)u_0+D(t)u_1$,
$u_t(t)=-A^{1/2}\sin(tA^{1/2})u_0+C(t)u_1$.
Dla $k\in\mathbb Z$ oznaczmy $\omega_k=\sqrt{1+k^2}$.
W zmiennych $(\omega_k\widehat u(k),\widehat v(k))$
operator przejścia $U(t)$ działa przez macierz
\[
 \begin{pmatrix}
 \cos(t\omega_k)&\sin(t\omega_k)\\
 -\sin(t\omega_k)&\cos(t\omega_k)
 \end{pmatrix}.
\]
Macierz ta jest unitarna, a wzory dodawania argumentów
dają własność grupy. Z tożsamości Parsevala
~\ref{thm:pde6-parseval} (z uwzględnieniem czynnika
$2\pi$ dla nieznormalizowanej miary) wynika,
że $U(t)$ jest izometrią na $Y$. Silna ciągłość
wynika najpierw dla skończonych sum Fouriera,
a następnie dla wszystkich elementów $Y$ przez
gęstość i jednostajne ograniczenie normy operatorów.

\textbf{Krok 2: lokalna lipschitzowskość.} Wymiar jeden
zapewnia zanurzenie $H^1(\T^1)\hookrightarrow C(\T^1)$.
Można je otrzymać bezpośrednio z nierówności
Cauchy'ego--Schwarza:
\[
 \sum_{k\in\mathbb Z}|\widehat w(k)|
 \le\left(\sum_{k\in\mathbb Z}\frac1{1+k^2}\right)^{1/2}
 \left(\sum_{k\in\mathbb Z}(1+k^2)|\widehat w(k)|^2\right)^{1/2}.
\]
Pierwszy szereg jest zbieżny, a drugi odpowiada
normie $H^1$ przez Parsevala. Szereg Fouriera
zbiega więc jednostajnie do ciągłego reprezentanta
$w$ i $\|w\|_\infty\le C\|w\|_{H^1}$.
Jeśli $\|v\|_{H^1},\|w\|_{H^1}\le R$, to
$\|g(v)-g(w)\|_2\le L_{CR}\|v-w\|_2
\le L_{CR}\|v-w\|_{H^1}$, gdzie
$L_{CR}=\max_{|s|\le CR}|g'(s)|$.

\textbf{Krok 3: punkt stały.} Dla $z=(u,v)$ definiujemy
\[
 \Psi(z)(t)=U(t)(u_0,u_1)
 +\int_0^tU(t-s)(0,g(u(s)))\,\mathrm d\lambda_1(s).
\]
Niech $R>\|(u_0,u_1)\|_Y+1$ i
$M_R=\sqrt{2\pi}\max_{|s|\le CR}|g(s)|$.
Silna ciągłość grupy i ciągłość superpozycji
$u\mapsto g(u)$ z $H^1$ do $L^2$ zapewniają
istnienie całki Bochnera i ciągłość $\Psi(z)$.
Na kuli $\|z\|_{C([0,T];Y)}\le R$ mamy
$\|\Psi(z)\|_{C([0,T];Y)}\le\|(u_0,u_1)\|_Y+TM_R$,
ponieważ $\|U(t-s)(0,h)\|_Y=\|h\|_2$.
Ponadto
\[
 \|\Psi(z)-\Psi(w)\|_{C([0,T];Y)}
 \le TL_{CR}\|z-w\|_{C([0,T];Y)}.
\]
Wybieramy $T$ tak, by $TM_R\le R-\|(u_0,u_1)\|_Y$
oraz $TL_{CR}<1$. Zasada Banacha daje punkt stały.
Pierwsza składowa wzoru ma postać
\[
 u(t)=C(t)u_0+D(t)u_1+
 \int_0^tD(t-s)g(u(s))\,\mathrm d\lambda_1(s).
\]
Mnożniki Fouriera pokazują, że
$D'(t)=C(t)$ w sensie silnym na $L^2$
oraz $C'(t)u_0=-A^{1/2}\sin(tA^{1/2})u_0$ w $L^2$.
Ponieważ $D(0)=0$, różniczkowanie całki w $L^2$
daje dokładnie drugą składową punktu stałego.
Jest ona ciągła w $L^2$. Równanie falowe zachodzi
w sensie dystrybucyjnym: każdy współczynnik Fouriera
spełnia wymuszone równanie skalarne
$\widehat u_k''+\omega_k^2\widehat u_k=\widehat{g(u)}_k$.
Przejście od skończonych sum testów Fouriera do
dowolnych testów gładkich jest uzasadnione ich
zbieżnością wraz z pochodnymi.
Dla dwóch rozwiązań o wspólnym ograniczeniu $R$
różnica wzorów całkowych i lemat Gronwalla dają
\[
 \|z(t)-\widetilde z(t)\|_Y
 \le e^{L_{CR}t}\|z(0)-\widetilde z(0)\|_Y .
\]
Dowodzi to jednoznaczności również poza pierwotną
kulą i lokalnej lipschitzowskiej zależności od danych;
czas istnienia jest wspólny dla danych z ustalonej
ograniczonej kuli.

Własność grupy $U(t)$ pozwala rozdzielić całkę
w dowolnej chwili $r$ i zapisać ten sam wzór
z danymi $z(r)$ oraz czasem liczonym od $r$.
Odwrotnie, rozwiązanie skonstruowane od tych danych
skleja się z poprzednim, gdyż po podstawieniu
wzoru na $z(r)$ spełnia pierwotne równanie całkowe.
Jednoznaczność na wspólnych przedziałach umożliwia
sklejenie rozwiązań do supremum czasów istnienia.
Gdyby dla skończonego $T_{\max}$ istniał ciąg
$t_n\uparrow T_{\max}$ z $\|z(t_n)\|_Y\le B$,
konstrukcja z ustaloną kulą o promieniu $B+2$
dawałaby rozwiązanie od każdej chwili $t_n$
na przedziale o tej samej dodatniej długości.
Dla dużego $n$ byłoby to przedłużenie poza
$T_{\max}$, sprzeczność. Dowodzi to podanej granicy.
\end{proof}

\begin{stwierdzenie}[Energia półliniowego równania falowego]
\label{prop:pde25-waveenergy}
Niech $g\in C^1(\R)$, $G(s)=\int_0^s g(r)\,\mathrm d\lambda_1(r)$,
a $(u,u_t)\in C([0,T];Y)$ będzie rzeczywistym
rozwiązaniem łagodnym z Twierdzenia~\ref{thm:pde25-wave}.
Wtedy
\[
 \mathcal E(t)=\int_{\T^1}\left(
 \tfrac12u_t^2+\tfrac12u_x^2+\tfrac12u^2-G(u)\right)
 \,\mathrm d\lambda_1(x)
\]
jest stała w czasie.
\end{stwierdzenie}
\begin{proof}
Dla rozwiązań o regularności
$u\in C([0,T];H^2)\cap C^1([0,T];H^1)
\cap C^2([0,T];L^2)$
mnożymy równanie przez $u_t$ i całkujemy po torusie.
Składnik $-u_{xx}u_t$ po całkowaniu przez części
jest pochodną $\tfrac12\|u_x\|_2^2$; pozostałe
składniki są pochodnymi $\tfrac12\|u_t\|_2^2$,
$\tfrac12\|u\|_2^2$ i
$-\int_{\T^1}G(u)\,\mathrm d\lambda_1$.
Suma pochodnych wynosi zero.

Dla ogólnych danych energetycznych nie zakładamy
istnienia $u_{tt}$ w $L^2$. Niech $P_N$ oznacza
rzut ortogonalny na mody Fouriera $|k|\le N$,
a $u_N=P_Nu$ i $v_N=P_Nu_t$.
Rzuty komutują z $A$; ze skalarnych równań
dla współczynników Fouriera otrzymujemy
\[
 u_N'=v_N,\qquad v_N'=-Au_N+P_Ng(u).
\]
Obie funkcje są klasy $C^1$ w skończenie wymiarowej
przestrzeni funkcji trygonometrycznych. Różniczkując
energię $\mathcal E_N$ zdefiniowaną tym samym
wzorem dla pary $(u_N,v_N)$, dostajemy
\[
 \mathcal E_N'(t)=
 \int_{\T^1}v_N\bigl(P_Ng(u)-g(u_N)\bigr)
 \,\mathrm d\lambda_1 .
\]
Nie jest to na ogół zero, ponieważ rzut rozwiązania
nie musi być rozwiązaniem równania nieliniowego.

Rzuty są kontrakcjami zarówno w $H^1$, jak i $L^2$
oraz zbiegają silnie do identyczności. Zbieżność jest
jednostajna na zwartych podzbiorach tych przestrzeni:
wystarczy przybliżyć zbiór skończoną siecią i użyć
jednostajnego ograniczenia norm rzutów.
Stosując to do zwartych obrazów $[0,T]$ przez
$u,u_t$ i $g(u)$, otrzymujemy jednostajną w czasie
zbieżność $(u_N,v_N)\to(u,u_t)$ w $Y$ oraz
$P_Ng(u)\to g(u)$ w $L^2$.
Z lokalnej lipschitzowskości $g:H^1\to L^2$
wynika również $g(u_N)\to g(u)$ jednostajnie w czasie
w $L^2$. Stąd
\[
 \sup_{0\le t\le T}|\mathcal E_N'(t)|
 \le \sup_t\|v_N(t)\|_2
 \sup_t\|P_Ng(u(t))-g(u_N(t))\|_2\longrightarrow0.
\]
Z zanurzenia $H^1\hookrightarrow C$ wynika jednostajna
zbieżność $u_N$ do $u$ także w zmiennych $(t,x)$.
Funkcja $G$ jest lipschitzowska na wspólnym ograniczonym
zakresie tych funkcji, więc $\mathcal E_N\to\mathcal E$
jednostajnie w czasie. Całkując tożsamość dla
$\mathcal E_N'$ od $0$ do $t$ i przechodząc do granicy,
otrzymujemy $\mathcal E(t)=\mathcal E(0)$.
\end{proof}

\begin{uwaga}[Energia nie zawsze kontroluje normę]
Jeżeli $G$ rośnie szybciej niż kwadratowo i z dodatnim
znakiem, składnik $-G(u)$ może sprawić, że energia
nie jest ograniczona od dołu. Jej zachowanie samo w sobie
nie daje globalnego istnienia. Dla nieliniowości
$g(u)=-u^3$ energia ma natomiast dodatni człon
$\int_{\T^1}u(t,x)^4/4\,\mathrm d\lambda_1(x)$
i kontroluje normę energetyczną.
\end{uwaga}

\begin{wniosek}[Globalność z koercywności energii]
\label{cor:pde25-wave-global}
Niech $g\in C^1(\R)$ i $G(s)=\int_0^s g(r)\,
\mathrm d\lambda_1(r)$. Jeśli dla pewnych
$0\le a<1/2$ i $b\ge0$ zachodzi
$G(s)\le a s^2+b$ dla każdego $s\in\R$,
to rozwiązanie łagodne z dowolnymi danymi w $Y$
istnieje dla wszystkich $t\ge0$.
W szczególności dotyczy to $g(s)=-s^3$.
\end{wniosek}
\begin{proof}
Niech $c=1/2-a>0$. Ze Stwierdzenia
~\ref{prop:pde25-waveenergy} oraz nierówności
dla $G$ wynika
\[
 \mathcal E(0)+2\pi b
 \ge \tfrac12\|u_t(t)\|_2^2+
 \tfrac12\|u_x(t)\|_2^2+c\|u(t)\|_2^2
 \ge c\|(u(t),u_t(t))\|_Y^2 .
\]
Użyliśmy $c\le1/2$ i tożsamości
$\|A^{1/2}u\|_2^2=\|u_x\|_2^2+\|u\|_2^2$
wynikającej z Parsevala.
Alternatywa eksplozji z Twierdzenia
~\ref{thm:pde25-wave} wyklucza skończony
maksymalny czas istnienia. Dla $g(s)=-s^3$
mamy $G(s)=-s^4/4\le0$, więc wystarczy $a=b=0$.
\end{proof}

\section{Eliptyczne równanie półliniowe i tożsamość Pohozaeva}
Rozważmy $-\Delta u=u^p$ w ograniczonym obszarze
$\Omega\subset\R^d$, $d\ge3$, z zerowym warunkiem
Dirichleta i $u>0$ wewnątrz. W przeciwieństwie do teorii
liniowej rozwiązanie może nie istnieć z przyczyn
wynikających z geometrii obszaru i wykładnika $p$.

\begin{lemat}[Tożsamość Pohozaeva]
\label{lem:pde25-pohozaev}
Niech $d\ge2$, $\Omega\subset\R^d$ będzie ograniczonym
obszarem klasy $C^2$, $F\in C^1(\R)$, $F(0)=0$,
a $u\in C^2(\overline\Omega)$ spełnia
$-\Delta u=F'(u)$ w $\Omega$ i $u=0$ na $\partial\Omega$.
Jeśli $\nu$ jest jednostkową normalną zewnętrzną, to
\begin{equation}\label{eq:pde25-pohozaev}
 \frac{d-2}{2}\int_\Omega|\nabla u|^2\,\mathrm d\lambda_d
 -d\int_\Omega F(u)\,\mathrm d\lambda_d
 =-\frac12\int_{\partial\Omega}(x\cdot\nu)
 (\partial_\nu u)^2\,\mathrm d\sigma(x),
\end{equation}
gdzie $\sigma$ jest miarą powierzchniową na brzegu.
\end{lemat}
\begin{proof}
Mnożymy równanie przez $x\cdot\nabla u$ i całkujemy.
Z jednej strony, ponieważ $\nabla F(u)=F'(u)\nabla u$,
$\int_\Omega F'(u)(x\cdot\nabla u)\,\mathrm d\lambda_d
=-d\int_\Omega F(u)\,\mathrm d\lambda_d$;
wyraz brzegowy $\int_{\partial\Omega}(x\cdot\nu)F(u)
\,\mathrm d\sigma$ znika z uwagi na $F(0)=0$.
Z drugiej strony całkowanie przez części daje
\[
 \begin{aligned}
 \int_\Omega(-\Delta u)(x\cdot\nabla u)\,\mathrm d\lambda_d
 &= -\int_{\partial\Omega}\partial_\nu u
 (x\cdot\nabla u)\,\mathrm d\sigma
 +\int_\Omega\nabla u\cdot\nabla(x\cdot\nabla u)
 \,\mathrm d\lambda_d\\
 &= -\int_{\partial\Omega}(x\cdot\nu)(\partial_\nu u)^2
 \,\mathrm d\sigma
 +\int_\Omega\bigl(|\nabla u|^2
 +\tfrac12 x\cdot\nabla|\nabla u|^2\bigr)
 \,\mathrm d\lambda_d.
 \end{aligned}
\]
Ponieważ $u$ ma stały ślad na brzegu, jego pochodne
styczne znikają tam, zatem $\nabla u=(\partial_\nu u)\nu$.
Twierdzenie o dywergencji zastosowane do
$x|\nabla u|^2$ zamienia ostatni składnik objętościowy
na $-\frac d2\int_\Omega|\nabla u|^2\,
\mathrm d\lambda_d+\frac12\int_{\partial\Omega}
(x\cdot\nu)(\partial_\nu u)^2\,\mathrm d\sigma$.
Po porównaniu z pierwszą stroną otrzymujemy
\eqref{eq:pde25-pohozaev}.
\end{proof}

\begin{twierdzenie}[Nieistnienie dodatnich rozwiązań przy wykładniku krytycznym]
\label{thm:pde25-pohozaev-nonexist}
Niech $d\ge3$, $p\ge(d+2)/(d-2)$ oraz
$\Omega\subset\R^d$ będzie ograniczonym obszarem klasy
$C^2$, \emph{ściśle gwiaździstym względem zera}:
$x\cdot\nu(x)>0$ dla każdego $x\in\partial\Omega$.
Nie istnieje funkcja $u\in C^2(\overline\Omega)$
taka, że $u>0$ w $\Omega$, $u=0$ na $\partial\Omega$
i $-\Delta u=u^p$ w $\Omega$.
\end{twierdzenie}
\begin{proof}
Mnożenie równania przez $u$ i całkowanie przez części
daje $\int_\Omega|\nabla u|^2\,\mathrm d\lambda_d
=\int_\Omega u^{p+1}\,\mathrm d\lambda_d$.
Stosujemy Lemat~\ref{lem:pde25-pohozaev}
do $F(s)=s^{p+1}/(p+1)$ dla $s\ge0$,
z dowolnym przedłużeniem $C^1$ na $\R$ i $F(0)=0$.
Otrzymujemy
\[
 \left(\frac{d-2}{2}-\frac{d}{p+1}\right)
 \int_\Omega|\nabla u|^2\,\mathrm d\lambda_d
 =-\frac12\int_{\partial\Omega}(x\cdot\nu)
 (\partial_\nu u)^2\,\mathrm d\sigma.
\]
Dla $p>(d+2)/(d-2)$ lewy współczynnik jest dodatni,
a prawa strona niedodatnia, więc $u\equiv0$,
sprzeczność. W przypadku krytycznym współczynnik
znika; ścisła gwiaździstość wymusza
$\partial_\nu u\equiv0$ na brzegu. Tymczasem $-\Delta u=u^p>0$
i $u>0$ wewnątrz, więc lemat Hopfa (Wniosek~\ref{cor:pde14-hopf-super})
dla obszaru klasy $C^2$ daje $\partial_\nu u<0$ na brzegu.
To sprzeczność kończąca dowód.
\end{proof}

\begin{uwaga}[Wykładnik krytyczny i ograniczenia twierdzenia]
Liczba $p_*=(d+2)/(d-2)$ odpowiada krytycznemu
wykładnikowi zanurzenia Sobolewa, ponieważ
$p_*+1=2d/(d-2)$, a to właśnie potęga $p+1$
występuje w całce z potencjału w tożsamości Pohozaeva.
Zanurzenie ma postać $H_0^1\hookrightarrow L^{2d/(d-2)}$.
Nie twierdzimy, że nie istnieją rozwiązania zmieniające znak
w przypadku krytycznym, ani że zachodzi ten sam wynik
na dowolnych obszarach: użyliśmy dodatniości oraz
ściśle gwiaździstej geometrii w twierdzeniu Pohozaeva.
\end{uwaga}

\section{Porównanie trzech mechanizmów i dalszy plan}
Dla równania ciepła z nieliniowością odpowiedni dobór
przestrzeni i kontrakcja dają lokalną jednoznaczność,
ale wzrost reakcji może wymusić eksplozję. Dla równania
falowego nieliniowość działa w układzie bez wygładzania,
co wymaga lokalnej lipschitzowskości nieliniowości
jako odwzorowania z $H^1$ do $L^2$. Dla zagadnienia
eliptycznego nawet bardzo gładkie dane i brzeg nie
zapewniają istnienia z powodu tożsamości całkowej.
W następnym rozdziale wykorzystamy operatory monotoniczne
i $p$-laplasjan, które prowadzą do innego mechanizmu istnienia
rozwiązań nieliniowych, bez kontrakcji Picarda.


% ROZDZIAL 26: monotonicznosc i plaplasjan
